\IntroVorlesung
\writeSectionFrame{\presentationTitle}{}

\begin{frame}{Vorgehensweise}
    \begin{itemize}
        \item Neues Thema mit vielen Beispielen
        \item Übungsblatt für Zuhause
        \item Besprechung und Vertiefung 
    \end{itemize}
    \vspace{0.3cm}
    \begin{center}
    \begin{tikzpicture}[
    >=Triangle,
    every node/.style={
        draw=Orange,
        fill=LightBlue,
        rounded corners,
        align=center,
        minimum width=2.6cm,
        minimum height=1cm
    }
]

\node (a) at (0,2) {Neues Thema};
\node (b) at (2.8,-1) {Übungsblatt};
\node (c) at (-2.8,-1) {Besprechung ÜB};

\draw[->, thick, Blue] (a) to[bend left=30] (b);
\draw[->, thick, Blue] (b) to[bend left=30] (c);
\draw[->, thick, Blue] (c) to[bend left=30] (a);

\end{tikzpicture}
    \end{center}
\end{frame}


\begin{frame}{Agenda}
\begin{columns}[c]

\column{0.40\textwidth}
\centering
\includegraphics<1->[width=\linewidth]{\BILDERPFAD/Database_Stock_Image.jpg}

\column{0.6\textwidth}
\begin{itemize}
    \item<2-> Datenbankentwurf
    \item<3-> Relationales Modell \& Algebra
    \item<4-> Structured Query Language (SQL)
    \item<5-> Relationale Entwurfstheorie
    \item<6-> Datenintegrität
    \item<7-> Transaktionsverwaltung
    \item<8-> Datenbanktuning
\end{itemize}

\end{columns}
\end{frame}